\textbf{Correction du QCM}
			
			\medskip
			Pour cette partie, une seule réponse est correcte pour chaque question.
			
			\medskip
			\textbf{\underline{Question 1}}
			
			\medskip
			On cherche la probabilité $P(\overline{A}\cap B)$.
			
			D'après l'arbre de probabilités, on lit :
			\[
			P(\overline{A})=0{,}6
			\qquad\text{et}\qquad
			P_{\overline{A}}(B)=0{,}7.
			\]
			
			On utilise la formule :
			\[
			P(\overline{A}\cap B)=P(\overline{A})\times P_{\overline{A}}(B).
			\]
			
			Ainsi :
			\begin{align*}
				P(\overline{A}\cap B)&=0{,}6\times 0{,}7\\
				&=0{,}42.
			\end{align*}
			
			\[
			\boxed{\text{Réponse b. }0{,}42}
			\]
			
			\medskip
			\textbf{\underline{Question 2}}
			
			\medskip
			On sait que $150$ élèves représentent $\dfrac35$ de l'ensemble des élèves de première générale.
			
			On note $N$ le nombre total d'élèves de première générale dans le lycée.
			
			On a donc :
			\[
			\frac35 N=150.
			\]
			
			Ainsi :
			\begin{align*}
				N&=150\times \frac53\\
				&=250.
			\end{align*}
			
			\[
			\boxed{\text{Réponse c. }250}
			\]
			
			\medskip
			\textbf{\underline{Question 3}}
			
			\medskip
			On considère :
			\[
			A=\frac13
			\qquad\text{et}\qquad
			B=\frac56.
			\]
			
			On calcule :
			\begin{align*}
				\frac{A}{B}+1
				&=\frac{\frac13}{\frac56}+1\\
				&=\frac13\times \frac65+1\\
				&=\frac25+1\\
				&=\frac25+\frac55\\
				&=\frac75.
			\end{align*}
			
			\[
			\boxed{\text{Réponse a. }\dfrac75}
			\]
			
			\medskip
			\textbf{\underline{Question 4}}
			
			\medskip
			La droite $d$ a pour équation réduite :
			\[
			y=\frac13x+1.
			\]
			
			Son coefficient directeur est $\dfrac13$. Si on se décale de 3 unités horizontalement, on remonte d'une unité verticalement pour retomber sur la droite.
			
			Son ordonnée à l'origine est $1$. Elle passe donc par le point de coordonnées $(0\ ;\ 1)$.
			
			La droite de la proposition c.
			
			\[
			\boxed{\text{Réponse c.}}
			\]
			
			\medskip
			\textbf{\underline{Question 5}}
			
			\medskip
			On développe :
			\[
			(x^3-1)^2.
			\]
			
			On utilise l'identité remarquable :
			\[
			(a-b)^2=a^2-2ab+b^2.
			\]
			
			Ici, $a=x^3$ et $b=1$. Donc :
			\begin{align*}
				(x^3-1)^2
				&=(x^3)^2-2\times x^3\times 1+1^2\\
				&=x^6-2x^3+1.
			\end{align*}
			
			\[
			\boxed{\text{Réponse b. }x^6-2x^3+1}
			\]
			
			\medskip
			\textbf{\underline{Question 6}}
			
			\medskip
			Une hausse de $20\ \%$ correspond au coefficient multiplicateur :
			\[
			1+\frac{20}{100}=1{,}2.
			\]
			
			Une baisse de $50\ \%$ correspond au coefficient multiplicateur :
			\[
			1-\frac{50}{100}=0{,}5.
			\]
			
			Le coefficient multiplicateur global est donc :
			\begin{align*}
				1{,}2\times 0{,}5&=0{,}6.
			\end{align*}
			
			La valeur finale représente donc $60\ \%$ de la valeur initiale.
			
			L'évolution globale est donc une baisse de :
			\[
			100\ \%-60\ \%=40\ \%.
			\]
			
			\[
			\boxed{\text{Réponse c. }40\ \%}
			\]
			
			\medskip
			\textbf{\underline{Question 7}}
			
			\medskip
			La classe compte $25$ élèves.
			
			Dans le tableau, on connaît déjà :
			\[
			8+7+4=19
			\]
			élèves.
			
			Il manque donc :
			\[
			25-19=6
			\]
			élèves dans la case correspondant aux élèves qui suivent la spécialité Mathématiques et qui ont plus de $16$ ans.
			
			Le tableau complété est donc :
			
			\begin{center}
				\begin{tblr}{colspec={Q[c,m]Q[c,m]Q[c,m]},hlines,vlines,width=0.7\linewidth,
						cell{1}{1}={white},cell{1}{2-3}={gray!45},cell{2-3}{1}={gray!45}}
					& \textbf{16 ans ou moins} & \textbf{Plus de 16 ans} \\
					\textbf{\shortstack{Suivent la spécialité\\ Mathématiques}} & 8 & 6 \\
					\textbf{\shortstack{Ne suivent pas la spécialité\\ Mathématiques}} & 7 & 4 \\
				\end{tblr}
			\end{center}
			
			Parmi les élèves de plus de $16$ ans, il y a :
			\[
			6+4=10
			\]
			élèves.
			
			Parmi ces $10$ élèves, $6$ suivent la spécialité Mathématiques.
			
			La probabilité cherchée est donc :
			\begin{align*}
				\frac{6}{10}&=\frac35.
			\end{align*}
			
			\[
			\boxed{\text{Réponse d. }\dfrac35}
			\]
			
			\medskip
			\textbf{\underline{Question 8}}
			
			\medskip
			On sait que $x$ et $y$ sont deux réels strictement positifs tels que :
			\[
			x=\frac{5}{2+y}.
			\]
			
			On multiplie les deux membres par $2+y$ :
			\[
			x(2+y)=5.
			\]
			
			On développe :
			\[
			2x+xy=5.
			\]
			
			Donc :
			\[
			xy=5-2x.
			\]
			
			Comme $x$ est strictement positif, on peut diviser par $x$ :
			\begin{align*}
				y&=\frac{5-2x}{x}\\
				&=\frac5x-\frac{2x}{x}\\
				&=\frac5x-2.
			\end{align*}
			
			\[
			\boxed{\text{Réponse d. }y=\dfrac5x-2}
			\]
			
			
			
			
			
\clearpage
